\chapter{Febrile Seizures}

\section*{Definition and Overview}
A febrile seizure is a convulsion triggered by fever in a young child, usually between 6 months and 5 years of age, who has no underlying neurological condition, infection of the brain, or metabolic disorder. Febrile seizures are the most common type of seizure in childhood, affecting around 3--5\% of children. They look alarming but are almost always harmless: the vast majority stop on their own within a few minutes, cause no brain damage, and do not mean the child has epilepsy. Simple febrile seizures, lasting under 15 minutes and involving the whole body, are by far the most common. This chapter explains what to do during a febrile seizure, when to seek emergency care, and how to manage the fevers that follow.

\section*{Epidemiology}
Febrile seizures affect 3--5\% of children, making them very common in Australian emergency departments and general practice. They are most frequent between 12 and 18 months of age, and boys are slightly more affected than girls. The risk is higher when a parent or sibling has had febrile seizures. Most children have a single episode, but about one in three has a recurrence, usually within the following year, particularly if the first seizure occurred at a young age or the fever was relatively low. After a simple febrile seizure, the risk of later epilepsy is only slightly higher than in the general population.

\section*{Aetiology and Pathophysiology}
\begin{itemize}
    \item \textbf{Fever:} any illness causing fever can trigger a seizure in a susceptible child; common triggers include viral infections, ear infections, and immunisation-related fever
    \item \textbf{Age-dependent susceptibility:} the developing brain between 6 months and 5 years is more excitable and prone to seizure activity, and this vulnerability fades with maturation
    \item \textbf{Genetic predisposition:} a family history of febrile seizures is common, and several genes influencing neuronal excitability have been identified
    \item \textbf{Infection type:} some viruses, such as human herpesvirus 6 (roseola) and influenza, are more strongly associated with febrile seizures
    \item \textbf{Mechanism:} the rapidly rising fever lowers the seizure threshold in an immature brain, allowing a brief generalised electrical discharge
    \item \textbf{Important reassurance:} the seizure is caused by the fever itself, not by a brain infection such as meningitis, which is always considered and excluded
\end{itemize}

\section*{Clinical Features}
\begin{itemize}
    \item A generalised convulsion: stiffening of the body followed by rhythmic jerking of the arms and legs
    \item Loss of consciousness, rolling eyes, and sometimes a brief cry at the start
    \item Duration of usually 1--3 minutes, and almost always less than 15 minutes (simple febrile seizure)
    \item Confusion, drowsiness, and sometimes irritability after the seizure (post-ictal state)
    \item Fever is present, which may be apparent before or after the seizure
    \item In complex febrile seizures: duration over 15 minutes, seizures confined to one side of the body, or more than one seizure within 24 hours; these need more thorough assessment
    \item A febrile seizure is not associated with a permanent change in the child's behaviour or development
\end{itemize}

\section*{Differential Diagnosis}
The key task is to distinguish a febrile seizure from more serious conditions.
\begin{itemize}
    \item Meningitis or encephalitis: infection of the brain or its coverings, suggested by a persistently unwell child, neck stiffness, and failure to settle
    \item Febrile status epilepticus: a febrile seizure lasting more than 30 minutes
    \item A first seizure of epilepsy, which can be fever-associated but recurs without fever later
    \item Rigors (shivering) with fever, which can look like a seizure but the child remains aware
    \item Breath-holding spells and syncope in the setting of a fever
    \item Other causes of acute neurological change, including low blood sugar or electrolyte disturbance
\end{itemize}

\section*{When to Seek Medical Care}
Any child who has a first febrile seizure should be assessed by a doctor to confirm the diagnosis and exclude a serious infection. Children who are unwell between fevers, or who have features of meningitis, need urgent review.

\textbf{Call 000 (or local emergency services) for:}
\begin{itemize}
    \item A seizure lasting more than 5 minutes
    \item The child not recovering, not breathing normally, or turning blue
    \item A seizure in a child under 6 months of age
    \item Repeated seizures without recovery between them
    \item Signs of meningitis: neck stiffness, severe headache, a rash that does not fade under pressure, or an unusually floppy, irritable, or drowsy child
\end{itemize}

\section*{Diagnosis and Investigations}
\begin{itemize}
    \item \textbf{History and examination:} the diagnosis is clinical, based on the age, the presence of fever, the appearance of the seizure, and a well child who recovers
    \item \textbf{Identifying the cause of the fever:} examination of the ears, throat, and chest, with urine testing in young children
    \item \textbf{Lumbar puncture:} performed if meningitis is suspected, particularly in infants under 12 months or in children who do not improve
    \item \textbf{Blood tests:} only when a specific concern exists, such as suspected sepsis or dehydration
    \item \textbf{EEG and brain imaging:} not needed after a simple febrile seizure; reserved for complex seizures or atypical features
    \item \textbf{Observation:} most children can be managed at home after assessment by a doctor, provided the cause of the fever is clear and the child recovers fully
\end{itemize}

\section*{Management}
\begin{itemize}
    \item \textbf{During the seizure:} place the child on their side, protect the head, remove nearby objects, time the seizure, and do not put anything in the mouth or restrain the child
    \item \textbf{Call for help:} seek emergency care if the seizure lasts beyond 5 minutes or the child does not recover promptly
    \item \textbf{Treat the fever:} paracetamol or ibuprofen and light clothing make the child comfortable, although they do not prevent further seizures
    \item \textbf{Emergency medicines:} a small group of children with prolonged or frequent febrile seizures may be prescribed midazolam (buccal) or diazepam to use at home during a seizure lasting more than 5 minutes
    \item \textbf{Treat the underlying illness:} antibiotics for bacterial infections such as ear or urinary infections
    \item \textbf{No routine daily medicines:} antiseizure medicines are not recommended to prevent febrile seizures because they are generally harmless events and the side effects outweigh any benefit
    \item \textbf{Reassurance and education:} families need clear guidance about what to do and what to expect
\end{itemize}

\section*{Complications}
\begin{itemize}
    \item Injury during the seizure, such as a fall or head bump, which is why supervision and a safe environment matter
    \item Prolonged seizure (febrile status epilepticus), the main rare serious risk
    \item Anxiety and distress in parents, which is common and understandable
    \item A slightly increased risk of later epilepsy, which remains small after a simple febrile seizure but is higher with complex seizures or a family history
    \item Rare complications of the underlying illness causing the fever
\end{itemize}

\section*{Prognosis and Outcomes}
The prognosis for febrile seizures is excellent. The seizure itself causes no brain damage, and children with febrile seizures have the same intellectual and developmental outcomes as other children. Most children outgrow them by age 5--6 as the brain matures. Around a third have further febrile seizures, but the frequency falls with age. The risk of epilepsy by adulthood is only about 1--2\% after a simple febrile seizure, compared with roughly 1\% in the general population, and is higher mainly after complex seizures or when there is a family history of epilepsy. Families can be confidently reassured that their child has experienced a frightening but benign event.

\section*{Prevention and Self-Care}
\begin{itemize}
    \item Learn first aid for seizures and make sure all carers know what to do
    \item Have a plan: know when to call 000 (seizure over 5 minutes) and when to see a doctor
    \item Give paracetamol or ibuprofen when the child has a fever and seems uncomfortable
    \item Dress the child lightly and encourage fluids to manage fever
    \item Do not give routine antiseizure medicines to prevent febrile seizures
    \item If the doctor prescribes emergency midazolam, learn how to give it and store it correctly
    \item Trust that the child will be well; speak to the doctor about any ongoing worry
\end{itemize}

\section*{Sources and Further Reading}
The following reputable sources provide further information for healthcare students and patients seeking reliable, current advice about febrile seizures.
\begin{itemize}
    \item Healthdirect Australia. \textit{Febrile seizures.} https://www.healthdirect.gov.au/febrile-seizures
    \item The Royal Children's Hospital Melbourne. \textit{Febrile convulsions fact sheet.} https://www.rch.org.au/
    \item Raising Children Network. \textit{Febrile convulsions.} https://raisingchildren.net.au/
    \item NHS. \textit{Febrile seizures.} https://www.nhs.uk/conditions/febrile-seizures/
    \item Mayo Clinic. \textit{Febrile seizure.} https://www.mayoclinic.org/
    \item MedlinePlus. \textit{Febrile seizures.} https://medlineplus.gov/
\end{itemize}
